% !TeX program = lualatex
% Standalone notebook; compile twice: lualatex 36.tex
% Research batch 39 down to 25, 27 September 2026.
% Original section preserved verbatim except for marked insertion blocks.
\documentclass[11pt,a4paper]{article}
\usepackage[margin=24mm,headheight=14pt]{geometry}
\usepackage{fontspec}
\setmainfont{Latin Modern Roman}
\setsansfont{Latin Modern Sans}
\setmonofont{Latin Modern Mono}
\usepackage{amsmath,amssymb,mathtools}
\usepackage{unicode-math}
\setmathfont{Latin Modern Math}
\usepackage{microtype}
\usepackage{enumitem,xurl,needspace}
\usepackage[dvipsnames]{xcolor}
\usepackage{hyperref}
\hypersetup{unicode=true,colorlinks=true,linkcolor=MidnightBlue,urlcolor=MidnightBlue,
 pdftitle={Research notebook: Problem 36},
 pdfauthor={Open Problems Math research notebook},bookmarksnumbered=true}
\usepackage{bookmark,titlesec,fancyhdr}
\titleformat{\section}{\Large\bfseries\raggedright}{\thesection}{0.7em}{}
\pagestyle{fancy}\fancyhf{}
\fancyhead[L]{\small\sffamily Open Problems Math\enspace|\enspace Research notebook}
\fancyhead[R]{\small\sffamily Problem 36}
\fancyfoot[L]{\small\sffamily 27 September 2026}
\fancyfoot[R]{\small\sffamily \thepage}
\setlength{\parindent}{0pt}
\setlength{\parskip}{5pt plus 1pt}
\setlength{\emergencystretch}{3em}
\setcounter{secnumdepth}{1}
\setcounter{tocdepth}{1}
\setlist[itemize]{leftmargin=3em,itemsep=5pt,topsep=4pt,parsep=0pt}
\allowdisplaybreaks
\newcommand{\N}{\mathbb{N}}
\newcommand{\Z}{\mathbb{Z}}
\newcommand{\Q}{\mathbb{Q}}
\newcommand{\R}{\mathbb{R}}
\newcommand{\C}{\mathbb{C}}
\newcommand{\F}{\mathbb{F}}
\newcommand{\A}{\mathbb{A}}
\newcommand{\PP}{\mathbb P}
\newcommand{\E}{\mathbb{E}}
\newcommand{\eps}{\varepsilon}
\newcommand{\dd}{\,\mathrm d}
\newcommand{\sref}[2]{\hyperlink{source:#1:#2}{[S#2]}}
\newcommand{\eref}[2]{\hyperlink{extra:#1:#2}{[E#2]}}
\newif\ifshowcatalogue
\showcataloguetrue
\newcommand{\cataloguescope}[1]{\ifshowcatalogue
 \paragraph{Exact scope in the supplied catalogue.}
 \begin{quote}\small #1\end{quote}\fi}
\numberwithin{equation}{section}
\begin{document}
\setcounter{section}{35}
% BEGIN PRESERVED SECTION
% ================================================================
% problemId: problem.extension-of-kroneckers-theorem-on-abelian-fields-to-any-algebraic-base-field
% source releaseRank: 36; source releaseStatus: open_with_solved_subcases
% exactTarget SHA-256: e83c72d0aa2f953d872f7859ed98b2e151e70e02a6ae890280d5db9039123c7f
\clearpage
\section{\texorpdfstring{Extension of Kroneckers Theorem on Abelian Fields to Any Algebraic Base Field}{Extension of Kroneckers Theorem on Abelian Fields to Any Algebraic Base Field}}\label{36}
\subsection{Definitions and mathematical statement}
A finite extension $L/K$ is abelian if it is Galois and $\operatorname{Gal}(L/K)$ is commutative. Hilbert's twelfth problem seeks an explicit function-theoretic construction of such extensions for every number field $K$, schematically
\[
 L=K\bigl(f_1(a_1),\ldots,f_r(a_r)\bigr),
\]
where the functions and allowed evaluation points are themselves explicitly specified from arithmetic data attached to $K$. Merely naming elements already known to generate $L$ is not such a construction. The catalogue records a broad construction program, not a uniquely specified function class or one fixed quantified theorem; admissible notions of explicitness must be stated in any proposed resolution.
\par\smallskip\textit{Formulation and concept references:} \sref{36}{1}.
\cataloguescope{For each number field K, construct generators for its finite abelian extensions as special values of explicitly defined analytic or arithmetic functions associated with K, extending roots of unity over Q and complex multiplication over imaginary quadratic fields.}
\subsection{Short English statement}
Can all abelian number-field extensions be generated by explicit special values, generalizing roots of unity and complex multiplication?
% BEGIN RESEARCH 36
\subsection{Research attempt}

\paragraph{Current frontier.}
This entry records a construction program, not a uniquely fixed allowable
function class. I retain its requirement that functions and evaluation points
be defined from arithmetic data, not from a generator already known to lie in
the desired extension. The ICMS formulation source discusses this program.
\sref{36}{1}. Dasgupta--Kakde prove a substantial totally real result:
Brumer--Stark units, together with specified square roots, generate the maximal
abelian extension, with a construction involving $p$-adic integration for
infinitely many primes. Their theorem is not restricted to imaginary quadratic
CM theory, but it also is not a theorem for arbitrary base number fields.
\sref{36}{2}, \eref{36}{1}. Honnor resolves a root-of-unity ambiguity in one
of the unit formulas; Dasgupta--Honnor--Spie\ss\ compare three formulas in a
2025 paper. \eref{36}{2}, \eref{36}{3}. The abstracts and theorem descriptions
were checked; no independent reproof of those arithmetic results is claimed.

\paragraph{Attack route.}
One possible analytic route generalizes reciprocity laws beyond the signatures
covered by existing special-value constructions. Its missing lemma has three
parts: algebraicity, the exact class-field action, and separation of distinct
Artin classes. A different route enumerates defining polynomials of abelian
extensions, but this does not meet the function-theoretic requirement. I
pursue the first route far enough to separate its genuinely arithmetic gap
from an avoidable primitive-element problem: once finitely many special
values separate Artin classes, an explicit bounded recipe produces generators
for every intermediate field.

\paragraph{Attempt.}
Let $H/K$ be finite abelian with group $G$, $N=|G|$, and suppose, for $r\ge1$, candidate
special values $u_1,\ldots,u_r$ have already been proved to belong to $H$.
Assume their Galois transforms can be described by a proposed reciprocity law
and that
\begin{equation}\label{r36:separation}
 \bigcap_{i=1}^r\operatorname{Stab}_G(u_i)=\{1\}.
\end{equation}
These are hypotheses, not consequences of numerical evaluation of the values.

\textbf{Lemma 36.1 (bounded primitive combination).}
Some integer $t$ with $0\le t\le(N-1)(r-1)$ makes
\[
 \alpha_t=\sum_{i=1}^r t^{i-1}u_i
\]
a generator of $H/K$. Here $t^0=1$, including at $t=0$.

\emph{Proof.}
For each $g\ne1$, the polynomial
$P_g(T)=\sum_i(g(u_i)-u_i)T^{i-1}$ in $H[T]$ is nonzero by
\eqref{r36:separation}. It has at most $r-1$ roots in $K$, hence at most that
many integer roots. At most $(N-1)(r-1)$ integers are bad for one or more
$g$. Among the displayed number of consecutive integers plus one, a good
one exists. Its stabilizer is trivial, so its orbit has $N$ elements. Since
$H/K$ is Galois, the degree of $\alpha_t$ over $K$ equals its orbit size and
$K(\alpha_t)=H$. The case $N=1$ is immediate, and $r=1$ is already covered
by the stabilizer condition. \hfill$\square$

The polynomial
\[
 Q_t(X)=\prod_{g\in G}(X-g\alpha_t)
\]
then lies in $K[X]$ by invariance, has degree $N$, and is irreducible. This
supplies an exact algebraic certificate once the reciprocity action and
inequalities $g\alpha_t\ne\alpha_t$ are established. It is more informative
than an unspecified generic linear combination: the finite range of possible
coefficients is given in advance.

\textbf{Lemma 36.2 (all intermediate fields from the same values).}
Let $\alpha$ have trivial stabilizer, let $J\le G$, and write
\[
 P_J(X)=\prod_{j\in J}(X-j\alpha)=X^{|J|}+c_1X^{|J|-1}+\cdots+c_{|J|}.
\]
Then $K(c_1,\ldots,c_{|J|})=H^J$.

\emph{Proof.}
The coefficients are fixed by $J$. If $g$ fixes all of them, $gP_J=P_J$;
since the roots $h\alpha$ are pairwise distinct, this says $gJ=J$, so
$g\in J$. Thus the common stabilizer of the coefficients is exactly $J$,
and Galois correspondence proves the equality. Because $G$ is abelian,
$H^J/K$ is Galois with group $G/J$. Lemma~36.1 can be applied again to the
coefficients to give one primitive generator with a bounded integer choice.
\hfill$\square$

This suggests a modular research program: construct a collection of values
for each ray modulus, prove that the reciprocity action separates the ray
classes, and let the two lemmas handle primitive generators and subextensions.
Finite sums and products of explicit values are again values of explicitly
specified finite combinations of the underlying functions, under a convention
allowing those arithmetic operations. That convention is stated here; it does
not permit inserting the unknown generator as a constant into a newly named
function.

An exact toy check is $H=\Q(\sqrt2,\sqrt3)$ with its four sign-changing
automorphisms. The two values $u_1=\sqrt2,u_2=\sqrt3$ separate the group.
The choice $t=1$, in the lemma's range $0\le t\le3$, yields
\[
 \prod_{\epsilon,\delta=\pm1}
 (X-\epsilon\sqrt2-\delta\sqrt3)=X^4-10X^2+1.
\]
The companion script checks this polynomial identity and its orbit size
exactly. This is a test of the finite mechanism, not a construction for a
new class field.

\paragraph{Stress test.}
The hard input is \eqref{r36:separation}, not the primitive-element theorem.
An invariant value gives no class-field information: if every $u_i$ lies in
a proper fixed field, every proposed combination remains there. Nonzero
values of an $L$-function or regulator do not automatically give algebraic
numbers with the correct individual Artin transforms. Algebraic independence
would also be the wrong objective for numbers that are supposed to lie in a
finite extension.

Approximations in $\C$ or $\C_p$ can sometimes certify that two values differ,
but finite precision does not prove algebraicity or exact field membership.
A recognition algorithm without a proved degree/height bound is not a proof.
Likewise formulas valid only up to unknown roots of unity can alter a
stabilizer, so that ambiguity cannot be dropped before applying either lemma.
The cited work resolving a specific ambiguity does not resolve every possible
one over arbitrary $K$. \eref{36}{2}.

The construction does not identify every number field as totally real or
imaginary quadratic, nor assume that passing to a larger field and taking
norms preserves enough class information. Norms can send all separating
values into a proper subfield; Lemma~36.2 uses the complete orbit polynomial
precisely to avoid relying on its constant term alone.

\Needspace{8\baselineskip}
\paragraph{Outcome.}
\textbf{Outcome: Conditional reduction.}

A bounded primitive-generator recipe and an exact intermediate-field extraction
lemma have been proved. They clarify a sufficient special-value package,
without claiming historical novelty for these elementary algebraic facts.
What remains is the central Hilbert-twelfth task for arbitrary $K$: construct
appropriate values, prove their algebraicity and reciprocity, and show they
separate the required classes. No such universal analytic construction is
supplied here, and the broad original program is not declared solved.

% END RESEARCH 36

\subsection{Sources}
\begin{itemize}\raggedright
\item[\textnormal{[S1]}]\hypertarget{source:36:1}{} International Centre for Mathematical Sciences, Recent progress on Hilbert's 12th problem (2024).
\url{https://icms.ac.uk/archive/workshop/recent-progress-on-hilberts-12th-problem/}.
{\small\textit{Catalogue formulation source.}}
\item[\textnormal{[S2]}]\hypertarget{source:36:2}{} Brumer–Stark Units and Explicit Class Field Theory.
\url{https://services.math.duke.edu/~dasgupta/papers/Integral-Gross-Stark.pdf}.
{\small\textit{Catalogue research/status reference; not a proof endorsement.}}
% BEGIN ADDED REFERENCES 36
\item[\textnormal{[E1]}]\hypertarget{extra:36:1}{} Samit Dasgupta and Mahesh Kakde,
\emph{Brumer--Stark Units and Explicit Class Field Theory}, Duke Mathematical
Journal 173 (2024), 1477--1555; arXiv:2103.02516.
\url{https://arxiv.org/abs/2103.02516}.
\item[\textnormal{[E2]}]\hypertarget{extra:36:2}{} Matthew H. Honnor,
\emph{On the Root of Unity Ambiguity in a Formula for the Brumer--Stark Units},
Canadian Mathematical Bulletin, DOI 10.4153/S0008439523001005.
\url{https://arxiv.org/abs/2302.12332}.
\item[\textnormal{[E3]}]\hypertarget{extra:36:3}{} Samit Dasgupta, Matthew H. L. Honnor
and Michael Spie\ss, \emph{On the equality of three formulas for Brumer--Stark units},
Journal of the London Mathematical Society 112 (2025), e70296.
\url{https://doi.org/10.1112/jlms.70296}.
{\small\textit{Consulted 27 September 2026; the totally real scope is retained.}}

% END ADDED REFERENCES 36
\end{itemize}

% END PRESERVED SECTION
\end{document}
